\section{View Changes (Leader Fault Tolerance)}
To maintain the system's liveness when the primary node fails or acts maliciously, nodes transition through a succession of configurations called \textit{views}.

Backup nodes monitor the primary's behavior using timers (\textit{timeouts}). A backup initiates a timer upon receiving a request; if the timer expires without any progress in processing, the backup assumes the leader is defective and halts the execution of the normal protocol.

The node ceases to accept normal operation messages and broadcasts a \texttt{VIEW-CHANGE} message requesting a transition to the next view ($v+1$). To be valid, this message appends cryptographic proofs of its latest stable checkpoint as well as all operations that had already reached the locally prepared state.

The new leader, designated by the formula $v \bmod N$ (where $N$ is the total number of nodes), takes control after aggregating at least $2f$ of these messages. It then issues a \texttt{NEW-VIEW} message containing consolidated proofs to resume the interrupted operations, ensuring that no operation already confirmed to a client is ever lost, and using null operations (\textit{no-ops}) to fill any empty slots in the sequence numbers. 

To prevent abuses, the protocol implements forced rotation and exponential timeout backing off in case of repetitive, failed view changes.